\subsection{Square ratios on the fitted grid}
\label{subsec:fitted-square-ratios}

Passing between tuples changes the scalar of a grid evaluation. The ratios on its constituent squares measure that change before the coefficient matrix is assembled.

Fix a selected point, its coefficient field \(K\), and its
marked collection of tuples. The square with indices \(i,j\)
has top dimensions \(a>a'\), signatures \((a-p,p)\) and
\((q,a'-q)\), and \(e=(a-a')/2\). Put \(e_0=e+q-p\).
The grid uses the stable range 
\[
  2(p+q)+1<\min(e,a'),
\]
adapted split planes and lattices, and regular infinity
exponents with gaps exceeding \(a+a'\). The real matching
maps for a tuple and its reference are taken on the same
geometric data. The base CM extension is unramified at
finite places over its real subfield, and its total field
is abelian over \(\Q\). Use the standard additive character
and the specified Weil frames.

At each inert place the
residue characteristic is odd and the residue cardinality
exceeds three. The additive character is unramified there
and at each split place in the selected set \(P\) of local
frames. The two splitting characters \(\alpha,\beta\) restrict
on the real-field ideles to the quadratic character of
\(E/F\) to powers \(a',a\), respectively, and have the prescribed
infinity exponents. Their base and auxiliary ramification
sets \(S_0,Q\) are disjoint, with 
\[
  S_0\cup Q\subset P.
\]
Both characters are unramified off that union, and their quotient
is ramified above every place in it. The matching place
\(v_1\) lies outside the union and splits completely in
\(E\). The real and finite maps use these restrictions in the
two matched Weil models at each adjacent step.

At the base place
\(v_0\in S_0\) the prescribed conductor exponent at the chosen
prime of \(E\) above \(v_0\) is one. Each extra level avoids
\(v_0\), so the toric character for each tuple retains this
conductor. It acts nontrivially on local units there.
Indeed conductor exponent one means that the local character \(\xi\) satisfies
\[
  \xi|_{1+\mathfrak p}=1,\qquad \xi|_{{\cal O}^{\times}}\ne1 .
\]
 Every
idele norm power is trivial on those units. The boundary
nonvanishing theorem for Hecke characters therefore applies at
\(s=1\), and the Euler product gives nonvanishing for real
\(s>1\). The reflected integer arguments used below are in the
range of the completed functional equation and the nonvanishing
comparison. \label{cor:fitted-square-ratios-1}

Write \(c(t)\) for the normalized top scalar, \(b_0(t)\)
for the scalar of the terminal definite triangle, and
\(\Lambda_t\) for the Hecke function of the square.
\begin{proposition}[The top ratio from the matched square]\label{prop:fitted-square-ratio-identity}
For the matched tuple and reference square with the stated
stable range and level hypotheses, the ratio \(r_{t,ij}=c(t)/c(t_0)\)
belongs to \(K^\times\) and satisfies \eqref{eq:fitted-square-tower}.
In particular, for each indicated pair of indices,
\[
 r_{t,ij}=c(t)/c(t_0)\in K^\times,\qquad r_{t_0,ij}=1.
\]

\end{proposition}

\begin{proof}
\proofstep{Use the same third test at the bottom.}

At a positive integer \(s\) in the square, put
\(\rho_t(s)=\Lambda_t(s)/\Lambda_{t_0}(s)\). The fitted embedding
identifies this nonzero ratio with an element of \(K\). At the definite
triangle the forward and reverse integrals are rational sums evaluated
against the same third test. Choose the test with nonzero denominator.
Dividing the two quotients, first on the triangle and then between
tuples, gives \(\gamma_t=b_0(t)/b_0(t_0)\in K^\times\).

\proofstep{Compare the scalar before and after one plane is removed.}
At any stage take the scalar in the order
\[
 I_{V,t}(\Phi)=c_{\rm before}(t) I_{V',t}(\Phi')
\]
on matched tests. Use the same order to define \(c_{\rm after}(t)\)
on the smaller triangle. Write \(d=(\dim V-\dim V')/2\) for the
half difference \emph{before} the removal, and let
\(\nu_t=\alpha_t/\beta_t\). In adapted frames the finite factor in
the plane comparison is
\[
 F_t(z)=\nu_t(\mathfrak b)^{\sigma_1}\nu_t(\mathfrak c)^{\sigma_2}
    \prod_{v\in S_0\cup Q_t}\nu_{t,v}(\sigma_3)
                        \varepsilon_v(z,\nu_{t,v},\psi_v).
\]

Here \(\mathfrak b,\mathfrak c\) are the fractional ideals on the
adapted isotropic lines in the member and in the joint space,
\(\sigma_1,\sigma_2,\sigma_3\in\{1,-1\}\) are the signs for this
type of step, and the local character uses the chosen split embedding.
For conductor exponent \(n\), uniformizer \(\varpi_v\), and
\(q_v=\mathrm Nv\), the normalization is
\[
 \varepsilon_v(z,\nu,\psi_v)=q_v^{-nz}\nu(\varpi_v)^n
   \sum_{x\in(\mathcal O_v/\varpi_v^n)^\times}
          \nu(\tilde x)^{-1}\psi_v(\tilde x/\varpi_v^n),
\]

where \(\tilde x\) is a unit lift. At the unramified places the
plane integral supplies the Euler factors of \(\Lambda_t\).
The Euler factors omitted at \(S_0\) are one because \(\nu_t\)
is ramified above every place in \(S_0\).

Take the constant term along the parabolic stabilizing the removed
isotropic line in the joint space. For the relevant Levi character,
the zero vector contributes evaluation at zero, denoted by \(E\). For the nonzero vectors one first averages
over the compact subgroup and then integrates along the line with
multiplicative weight \(\nu_t(l)|l|_E^s\). Denote the resulting test
by \(J_s\). Splitting the line into its zero and nonzero vectors
places \(E\) on the unchanged member and \(J_s\) on the member
losing a plane.

In the lower matched quotient the pairs of tests are
\[
 \begin{array}{c|c|c}
 \text{reduced member}&\text{lower tests}&
                    \text{multiplier }M_t\\ \hline
 V' & (M_t E\Phi,\ J_{d+1}\Phi') & D F_t(-d)\Lambda_t(d+1)\\
 V  & (J_d\Phi,\ M_t E\Phi') & D F_t(1-d)\Lambda_t(d).
 \end{array}
\]
Here \(D\ne0\) is the common real and measure factor in each row.
The finite scalar is obtained by multiplying the local integrals.
At an unramified place the valuation shells give a geometric series.
At a ramified split place orthogonality on the unit group leaves the
shell of the conductor, with the character sum displayed above. The
two line ideals account for the changes of generators.

Write \(I_1,I'_1\) for the integrals on the lower triangle. In the
first row the lower matching identity and the upper constant-term
identity, evaluated on a test with nonzero lower value, read
\[
 \begin{aligned}
 M_t I_1(E\Phi)&=c_{\rm after}(t) I'_1(J_{d+1}\Phi'),\\
 D_1 I_1(E\Phi)&=c_{\rm before}(t)D_2 I'_1(J_{d+1}\Phi').
 \end{aligned}
\]
The constants \(D_1,D_2\ne0\) depend on the Haar measures and geometry
of this step. Eliminating the two integrals gives
\(c_{\rm before}(t)M_t=(D_1/D_2)c_{\rm after}(t)\).
For the second row the same comparison reads
\[
 \begin{aligned}
 I_1(J_d\Phi)&=c_{\rm after}(t) M_t I'_1(E\Phi'),\\
 D_1 I_1(J_d\Phi)&=c_{\rm before}(t) D_2 I'_1(E\Phi').
 \end{aligned}
\]
Here \(c_{\rm before}(t)=(D_1/D_2)M_t c_{\rm after}(t)\).

Absorbing the nonzero geometric constants into one \(C\) for each
step, we obtain
\[
 \begin{aligned}
 c_{\rm before}(t)\Lambda_t(d+1)
       &=[C F_t(-d)]^{-1} c_{\rm after}(t)
       &&(\dim V'\mapsto\dim V'-2),\\
 c_{\rm before}(t)
       &= C F_t(1-d)\Lambda_t(d)c_{\rm after}(t)
       &&(\dim V\mapsto\dim V-2).
 \end{aligned}
\]
In the second row all planes on the primed side have already
been removed. Within either step the same \(C\) occurs for
\(t,t_0\) because the real maps and adapted geometric data agree.

\proofstep{Telescope first on the smaller member and then on the larger.}
Let \(r_k=c_k(t)/c_k(t_0)\) after \(k\) removals, so that
\(r_0=r_{t,ij}\). At step \(k\) let \(u_{k+1}\) be the ratio of
the finite factors as they occur in the corresponding line above.
Thus in adapted frames the two recurrences are
\[
\begin{aligned}
 r_b\rho_t(e+b+1)&=u_{b+1}r_{b+1} &&(0\le b<q),\\
 r_{q+h}&=u_{q+h+1}\rho_t(e+q-h)r_{q+h+1} &&(0\le h<p).
\end{aligned}
\]

In the first line \(u=F_{t_0}(-d)/F_t(-d)\), with \(d=e+b\).
In the second \(u=F_t(1-d)/F_{t_0}(1-d)\), with \(d=e+q-h\).
Changing adapted finite frames multiplies these ratios by character
values, which are units away from the extra level. Successive
substitution gives the following two expressions involving \(r_q\).
\[
\begin{aligned}
 r_0\prod_{s=e+1}^{e+q}\rho_t(s)
   &=\Bigl(\prod_{k=1}^{q}u_k\Bigr)r_q,\\
 r_q&=\Bigl(\prod_{k=q+1}^{q+p}u_k\Bigr)
         \Bigl(\prod_{s=e_0+1}^{e+q}\rho_t(s)\Bigr)r_{p+q}.
\end{aligned}
\]

An empty product is one. On identifying the last triangle with the
chosen definite triangle, the common real scalar cancels in the ratio.
Changing the bottom lattice and frames gives a factor \(\zeta\in K^\times\)
which is a unit away from the extra level. Thus
\(r_{p+q}=\zeta\gamma_t\) and
\begin{equation}\label{eq:fitted-square-tower}
 r_{t,ij}\prod_{s=e+1}^{e+q}\rho_t(s)
     =\zeta\gamma_t\prod_{k=1}^{p+q}u_k
                        \prod_{s=e_0+1}^{e+q}\rho_t(s).
\end{equation}

The finite quotients are products of algebraic character values and
ratios of nonzero Gauss sums in the fitted field. Division by the
nonzero product on the left gives \(r_{t,ij}\in K^\times\).
For \(t=t_0\) the ratio is one. \label{cor:fitted-square-ratios-2}

\end{proof}

Fix a finite place \(w\) above the rational prime \(\ell\).
Set \(e_w=v_w(\ell)\) and \(t=t(w)\).

\begin{proposition}[The bottom norm and the reflected indices]\label{prop:fitted-square-bottom-index}
At \(w\mid\ell\), use \(t=t(w)\), with its extra level prime to
\(\ell\), or the base tuple. For \(s\le0\), the left side below uses the ratio normalized by
the completed functional equation; at the reflected positive index the ratio gives the same valuation directly. In the indicated range,
\[
 v_w(\rho_t(s))=v_w(\rho_t(1-s))\qquad(1-e_0\le s\le0,\ s\in\mathbb Z).
\]

The integral comparisons for the two bottom norms give
\eqref{eq:fitted-square-bottom}; its clearing cost includes both
\(v_w(A)\) and \(v_w(A')\).
\end{proposition}

\begin{proof}
In the nonbase case the
twist level is prime to \(\ell\). The factor \(\zeta\)
is then a unit at \(w\), and the quotient of local factors
\(u_h\) has absolute valuation at most
\[
 B_w^{\rm step}=2\sum_{v\in S_0} c_v v_w(\mathrm Nv).
\]
Indeed the base conductors are unchanged in the two
comparisons. At a base place with unchanged character on
units, the Gauss quotients cancel up to unit factors. At
the remaining base places the two nonzero integral Gauss
factors divide the same conductor norm power. The new local
factors have norms supported on the prime-to-\(\ell\)
twist level, so their quotients are units at \(w\).
\proofstep{Reflect the integer arguments.}

At \(1-e_0\le s\le0\) the functional equations express the
ratio as an away unit times \(\rho_t(1-s)\). Define
\(\rho_t(s)\in K^\times\) by this comparison at these indices.
Thus 
\[
  v_w(\rho_t(s))=v_w(\rho_t(1-s)).
\]

\proofstep{Compare the integral coordinates at each bottom norm.}

Let \(A,A'\) be the two base definite norm values.
The first-occurrence product for the first twisted norm is
\(Z_t=A\prod_{1-e_0\le s\le h_0}\rho_t(s)\), with
\(e_0\le h_0\) and \(h_0+p+q\) equal to the square half rank.
Both norms are nonzero. The integral coordinate comparisons
for the two bottom lifts express \(Z_t/b_1\) and \(b_2A'\)
by bottom mass sums, with 
\[
  b_1/b_2=\gamma_t
\]
 times an
away unit.

Choose a local integral representative in each
class. The nonzero terms in either mass sum have one reciprocal
stabilizer and two values in the integral coordinates. Their valuations
are bounded below by \(-3\delta_w\) since each of the three
terms has valuation at least \(-\delta_w\) and the comparison
matrix has integer entries. Taking the product of the two sums
and using the ultrametric inequality gives
\begin{equation}\label{eq:fitted-square-bottom}
 v_w(\gamma_t)\le v_w(A)+v_w(A')+
             \sum_{s=1-e_0}^{h_0} v_w(\rho_t(s))+6\delta_w .
\end{equation}
In particular this uses two coordinate identities with opposite
placements of the triangle constants. The unshifted base
norms also have \(v_w(A),v_w(A')\ge-3\delta_w\) from their
integral mass sums. \label{cor:fitted-square-ratios-3}

\end{proof}

\subsubsection{The toric shift and its branch}

The toric deletion has two possible branches. A depth shift determined on the common base polynomial permits both to be measured in the same coordinates.

\begin{proposition}[A shift from the common base polynomial]\label{prop:fitted-toric-denominator-shift}
For each positive square index \(s\), the base norm gives an
integer \(g_w\) with
\[
  -n_s(\ell)e_w\le v_w(\rho_t(s))+g_w .
\]
The shift is obtained from the base polynomial, independently of
the deletion branch; the two deletion scalars \(d_c\) use the same
period.
\end{proposition}

\begin{proof}
For each positive index \(s\) in the square, let \(T_0\) be
its base toric value, and let \(T_0^{(0)},T_0^{(1)}\) be the
values for the two deleted-character branches. The branch
\(c(w)\) satisfies that its character is the base character
multiplied by the grid's deleted character at \(\ell\).

Write \(d_c=T_0^{(c)}/T_0\in K^\times\).
The scalar in the base toric slot is \(1\).
Use the same period and normalization for all three twelfth
powers 
\[
  \beta_0,\beta_c=d_c^{12}\beta_0.
\]
 The toric
lattice formula in Section~\ref{subsec:toric-twist}
writes the base values with the original conductor and the
tuple values with extra split places disjoint from \(\ell\).
The tuple character differs by a ring class character of the
enlarged order. With \(n_s(\ell)\) the ceiling of a twelfth of
the class-value denominator allowance, this formula and its
integrality estimate give
\(\ell^{12n_s(\ell)}\rho_t(s)^{12}\beta_0\)
integral at every place above \(\ell\). The same clearing
works for each of \(\beta_0,\beta_c\) at the base order.

Let \(P_0\) be the monic minimal polynomial of \(\beta_0\)
over \(K\) in the fitted embedding. Define the real normalized
valuation \(m_w=v_w(P_0(0))/(12\deg P_0)\), and let
\(g_w\) be its integer truncation toward zero. The norm
from \(K(\beta_0)\) to \(K\) gives

\[
  -n_s(\ell)e_w\le v_w(\rho_t(s))+m_w.
\]
The two integer terms show that this remains true with \(g_w\).

\end{proof}

\subsubsection{The simultaneous toric comparison}

The separate toric bounds will be summed over places. Their shifts and the assigned tuples must therefore be held fixed when taking both sides of the estimate.

\begin{proposition}[Two-sided toric bounds for the assigned tuples]\label{prop:fitted-toric-assigned-bounds}
For the selected places and their assigned twists, the two inequalities \eqref{eq:fitted-toric-two-sides} at each positive square index combine to give
\[
 -a_s(w)-n_s(\ell)e_w\le v_w(\rho_t(s))\le a_s(w)+e_w/\ell^2.
\]
 The bounds for the shift and the two deletion branches are
\eqref{eq:assigned-weighted-deletions}, with \(B_T=Cfb(\log b+\log x)\).
These bounds use the same \(g_w\) from the base minimal polynomial.
\end{proposition}

\begin{proof}
\proofstep{Use the simultaneous comparison for every cyclic power.}

The marked coefficient places are grouped in conjugate pairs.
For a nonbase place \(w\), the chosen twist index \(i_w\) lies under
one of \(w,\bar w\), has residue characteristic \(\ell\), and
carries the marked split auxiliary level \(Q_w\).
On each selected pair the record carries the same index and
the character with the prescribed side convention. The
clearing exponent \(N_0(i)\) of Section~\ref{subsec:toric-twist}
is chosen for the abelian coefficient functions at index \(i\)
before the auxiliary primes and ray characters. There are
finitely many indices on the datum.

At a fixed index \(i\), write \(n_{i,\lambda}\) for the clearing exponent obtained for each abelian label with an allowed specialization. The finite maximum at this index is taken before selecting the auxiliary characters
\[
  N_0(i)=\max_\lambda n_{i,\lambda}.
\]
For an empty set of labels the clearing condition is vacuous and one may take zero.
Apply Section~\ref{subsec:retained-twist} with allowance
\(1/\ell^2\) on each pair, adding the rational primes already
used to the avoidance set. This supplies the level, the
coefficient extension, and a ray character \(\psi'_w\) of
\(\ell\)-power order. Write \(\ell^n\) for the order of \(\psi'_w\). For the second reading, inversion replaces \(a\) by \(-a\).

In particular every unit exponent in either reading generates the same character group
\[
  (a,\ell)=1\quad\Longrightarrow\quad \langle(\psi'_w)^{-a}\rangle=\langle(\psi'_w)^a\rangle=\langle\psi'_w\rangle,\qquad \operatorname{ord}((\psi'_w)^a)=\ell^n.
\]
For every marked \(w\), every \(a\) prime to \(\ell\), and every
toric label \(\lambda\) of its index, the signed Fourier coefficient \(T_{w,\lambda}(a)\) at \((\psi'_w)^a\) satisfies
\begin{equation}\label{eq:fitted-toric-all-powers}
 T_{w,\lambda}(a)\ne0,\qquad
 -N_0(i_w)\le v(T_{w,\lambda}(a))
             \le v(S_{i_w,\lambda})+\ell^{-2}.
\end{equation}
Here \(S_{i,\lambda}\) is the empty-level sum and \(v(\ell)=1\).
The same selected level and character satisfy the other
constraints of the simultaneous specialization. The proof uses
the unique lowest term at every primitive root of the selected
order, so these assertions hold for all coprime powers on this
one choice. The field fit carries this property into the
evaluation grid. \label{cor:fitted-square-ratios-4}

\proofstep{Read the deleted branch and the valuation side separately.}

In the square at the critical integer \(s\), the entry labels are transported by
\[
   (i,j,s)\longmapsto (m-i,j,s),
\]
with the harmonic exponents of that square at \(s\).
The prime deletion multiplies the base square character by
the deleted character at \(\ell\). The resulting branch
is \(c(w)\) above. Independently, the prime below \(w\) or
\(\bar w\) chooses the side of the valuation comparison.
Use the corresponding toric label \(\lambda\) on that side.
The class-sum comparison groups order classes by their inverse
Artin labels.
For the class weights \(a_x\) this regrouping has the form
\[
  \sum_x a_x=\sum_\lambda\ \sum_{\operatorname{Art}(x)^{-1}=\lambda}a_x.
\]
 At the empty level it identifies \(S_{i_w,\lambda}\)
with a common nonzero multiple of \(T_0^{(c(w))}\).

At the
selected level it identifies \(T_{w,\lambda}(a)\) with the
same multiple of the toric sum with character \(\psi_c^{ca}\),
where \(\psi_c\) is the ring class character of the complex
reading and \(c\) is prime to \(\ell\). The identification of
the exponents follows by identifying the two cyclic groups of
character values under the chosen coefficient embeddings. More explicitly, with \(A_w\ne0\) the common proportionality
factor in the matched readings, the empty branch and its
specialization read
\[
 S_{i_w,\lambda}=A_w T_0^{(c(w))},\qquad
 T_{w,\lambda}(a)=A_w T_{\psi_c^{ca}}^{(c(w))}
\]
after embedding in the common field. Here \(T_\eta^{(c(w))}\)
denotes the toric sum for the indicated branch and character.
The exponent \(c\) absorbs inversion on the second side.
The choice of branch \(c(w)\) is independent of that inversion.

\proofstep{Match the Galois centres and normalize the valuation.}

To compare at \(w\), extend the local coefficient embedding
to a common valued field. The finite extension in the field
fit is Galois over the smaller coefficient field. Transitivity
on primes above the common lower prime lets one choose a
field automorphism giving the required centre.
Write \(L/K_1\) for the Galois extension in this comparison. On the two chosen extensions \(w_1,w_2\) of the lower place, transitivity gives
\[
  \exists\,\sigma\in\operatorname{Gal}(L/K_1)\quad \sigma(w_1)=w_2.
\]

On the conjugate
side compose the comparison with complex conjugation. The
algebraic character values in the base class prefactors and
their conjugates are in the smaller field by the principal-idele
type and the finite class representatives. The CM embeddings
and the quadratic period normalization are those of the fitted
frame. With these identifications the toric Galois comparison changes the
period-normalized ratio by a root of unity and sends
\(\psi_c\) to \(\psi_c^r\) with \(r\) prime to \(\ell\).
In particular, for the \(\ell\)-power order \(n\) of \(\psi_c\),
\[
  (r,n)=1\quad\Longrightarrow\quad\operatorname{ord}(\psi_c^r)=\operatorname{ord}(\psi_c)=n .
\]
In fact it acts on the cyclic character values by an automorphism.
If the order is \(\ell^n\), choose an integer \(a\) prime to
\(\ell\) with \(ca\equiv r\pmod{\ell^n}\).

For the trivial
character take \(a=1\). This choice occurs inside the
all-powers assertion \eqref{eq:fitted-toric-all-powers}.
The nonzero empty-level scalar cancels in the ratio.
The root of unity has valuation zero, and two nonzero values of equal valuation in the local fraction field
differ by a unit. Thus in that field, with valuation ring
\(O_w^{\mathrm{loc}}\),
\[
   \frac{\rho_t(s)}{d_{c(w)}}S_{i_w,\lambda}
       =uT_{w,\lambda}(a),\qquad u\in (O_w^{\mathrm{loc}})^\times .
\]
This gives the upper inequality
\[
 v_w(\rho_t(s)/d_{c(w)})\le e_w/\ell^2 .
\]
Indeed extension of valuation multiplies additive
valuations by ramification and leaves their normalization
by \(v(\ell)\) unchanged. The unit has valuation zero.
Hence with
\(a_s(w)=|g_w|+|v_w(d_{c(w)})+g_w|\) we have
\begin{equation}\label{eq:fitted-toric-two-sides}
  -v_w(\rho_t(s))\le a_s(w)+n_s(\ell)e_w,\qquad
  v_w(\rho_t(s))\le a_s(w)+e_w/\ell^2.
\end{equation}
\label{cor:fitted-square-ratios-5}

\proofstep{Bound the shift and deletion terms by the base norm heights.}

For the finite sums of \(a_s\), the constant term of the
minimal polynomial of \(\beta_c\) gives
\(m_{c,w}=m_w+v_w(d_c)\). The product formula and the norm
comparison via \(K(\beta_c)\) express
\(\sum_w\omega_w m_{c,w}\) as the normalized logarithm of the
absolute constant coefficient of its minimal polynomial over
\(\Q\), divided by twelve. Use the same calculation for the
base \(m_w\).

If these archimedean means are at most \(A_T\)
and \(\sum_\ell n_s(\ell)\log\ell\le D_T\), then the weighted
sum of the absolute values is at most \(A_T+2D_T\) for each
of the three normalized valuations. Indeed, for each of these three valuations \(m'_w\), write
\((m'_w)^-=\max(0,-m'_w)\). The local bound gives
\[
 \sum_w\omega_w |m'_w|
 =\sum_w\omega_w m'_w+2\sum_w\omega_w (m'_w)^-
 \le A_T+2\sum_\ell n_s(\ell)\sum_{w\mid\ell}\omega_w e_w
 \le A_T+2D_T .
\]
Here the full fibre has weight \(\sum_{w\mid\ell}\omega_w e_w=\log\ell\).
\proofstep{Average over the quadratic ideal classes.}

For the archimedean comparison, the CM transport writes
each conjugate normalized twelfth power as the twelfth
power of a Hecke value at permuted weights times the
period factor. For a conjugating ideal \(\mathfrak a\) in
the imaginary quadratic subfield of discriminant \(-d\),
the logarithm is bounded by
\[
 12\log C_0+12\log |L_{\mathfrak a}(s)|
       -M\log((\mathrm N\mathfrak a)^6|\Delta_{\mathfrak a}|),
 \qquad M=2sf+2\sum_\tau\kappa_\tau .
\]
Here \(\Delta_{\mathfrak a}\) is the period lattice
discriminant and \(C_0\) is the CM normalization constant.
The logarithm of the norm-discriminant factor depends only
on the ideal class. Over embeddings fixing the quadratic
subfield the classes occur equally often, and conjugating
exchanges the two quadratic embeddings without changing
absolute values. The Chowla--Selberg mean of that
logarithm is
\[
 12\log(2\pi)-6\log d+6\gamma-6L'(1,\chi_d)/L(1,\chi_d),
\]
where \(\chi_d\) is the quadratic character of discriminant
\(-d\). We use the Chowla--Selberg formula in the same normalization
as Lemma~\ref{lem:flat-cm-degree}. \cite{colmez-chowla-selberg}

Here are the size estimates needed to take this mean uniformly. Put
\(L_b=\log b\), \(f=[F:\Q]\), \(K_\kappa=\sum_\tau\kappa_\tau\),
and \(X_s=2sf\). The fitted square uses the field and size conditions
of Lemma~\ref{lem:window-parameter-choice} and
Proposition~\ref{prop:translate-completion}. In particular its ranks,
weights and field satisfy
\[
 16\le X_s\le b,\quad 1\le\kappa_\tau\le5b,\quad
 \log |D_F|\le fL_b,\quad
 \log d+\left|\frac{L'}{L}(1,\chi_d)\right|\ll L_b .
\]
For the last estimate only an upper bound for the absolute logarithmic
derivative of this selected quadratic character is needed. The
normalization constant in the toric class formula has the bound
\[
 \log C_0\le\tfrac{2s-1}{2}\log|D_F|
      +\sum_\tau\log\Gamma(2s+\kappa_\tau)+8K_\kappa
      \ll fbL_b.
\]
Indeed \(2s+\kappa_\tau\le7b\) and
\(\log\Gamma(n)=\log((n-1)!)\le n\log n\).
The CM period term in the mean has the same order because
\(M=X_s+2K_\kappa\ll fb\).

For clarity, let \(\eta\) denote any of the conjugate unitary Hecke
characters in this calculation, with conductor \(\mathfrak f_\eta\).
Write \(Z=\zeta(3/2)^{2f}\) and
\(A=|D_E|\mathrm N\mathfrak f_\eta\).
Use the primitive \(L\)-function of \(\eta\) in the functional equation.
If Euler factors above the fixed set \(S_*\) are omitted in the toric
value, each inverse local factor has modulus at most \(2\) for
\(\RePart z\ge1\). The omitted factors have total cost bounded by
\[
 \log\prod_{v\mid S_*}|1-\eta(v)(\mathrm Nv)^{-z}|
      \le 2f |S_*|\log2\le10f\log2 .
\]

For \(\RePart z\ge3/2\) the Euler product gives \(|L(z,\eta)|\le Z\),
in particular at every integer \(s\ge2\).
At \(s=1\) put \(m_\tau=1+\kappa_{\pi(\tau)}\), where \(\pi\) permutes the
CM embeddings, and write the archimedean factor as
\[
 G(z)=\prod_\tau 2(2\pi)^{-z-m_\tau}\Gamma(z+m_\tau).
\]

The integral for the gamma function gives
\(|G(3/2+it)|\le G(3/2)\). Thus \(|G(z)L(z,\eta)|\le G(3/2)Z\)
on the right edge of \(-1/2\le\RePart z\le3/2\).
The functional equation for \(\eta\) and its inverse gives the bound
\(A G(3/2)Z\) on the left edge. The nonzero infinity type makes the
completed function entire, and its boundedness in vertical strips
permits the Phragm\'en--Lindel\"of estimate at \(z=1\).
These analytic properties are the global zeta-integral theorem.
\cite[Sections~4.4--4.5]{tate-hecke}
Since \(\Gamma(3/2+m)\le\Gamma(2+m)=(1+m)\Gamma(1+m)\) for \(m\ge1\),
division by \(G(1)\) gives
\[
 |L(1,\eta)|\le ZA\prod_\tau(1+m_\tau).
\]

The base conductor has exponent at most two at the places above
\(S_*=\{g_0,g_1,p_d,p'_d,q_0\}\) and exponent zero elsewhere.
The size bounds for the level primes give
\(\max_{\ell\in S_*}\log\ell\ll L_b\).
Consequently
\[
 \log\mathrm N\mathfrak f_\eta\le4f\sum_{\ell\in S_*}\log\ell\ll fL_b,
 \qquad \log|D_E|=2f\log\operatorname{rd}_E\ll fL_b .
\]

The bound for \(L(1,\eta)\) is uniform in the permutation and in
the two deletion branches. Together these estimates bound the
archimedean mean divided by twelve by \(A_T\ll fbL_b\).

\proofstep{Sum the local denominator exponents.}
In Proposition~\ref{prop:toric-class-denominator} write \(c_E(\ell)\)
for the exponent contributed by the partial zeta values of \(F\).
For every ordinary ideal class \(C\), put
\(z_C=\zeta_F(1-2s,C)\).
Deligne and Ribet's integrality theorem applied to the indicator of
\(C\) and the principal ideal \((a)\), for an integer \(a>1\)
prime to \(\ell\), gives
\[
              (a^{X_s}-1)z_C\in\mathbb Z_\ell .
\]
To see the passage from their imprimitive values to \(z_C\), omit
the primes over \(\ell\) first. Restoring the factor for
\(\mathfrak p\mid\ell\) in the vector indexed by ideal classes
inverts \(1-(\mathrm N\mathfrak p)^{2s-1}P_{\mathfrak p}\),
where \(P_{\mathfrak p}\) translates by the inverse class of
\(\mathfrak p\). This integral matrix is the identity modulo \(\ell\),
so its inverse preserves the same integral lattice.
\cite[Theorem~0.4]{deligne-ribet-values}

At odd \(\ell\), choose \(a\) generating the units modulo \(\ell^2\).
If \(\ell-1\nmid X_s\), then \(a^{X_s}-1\) is a unit. Otherwise
its valuation is \(1+v_\ell(X_s)\), by binomial expansion. At two
take \(a=3\). Since \(X_s\) is even,
\(v_2(3^{X_s}-1)=2+v_2(X_s)\). Dividing the partial zeta values
by \(2^f\) as in the definition of \(c_E\) therefore gives
\[
 \begin{split}
 c_E(2)&\le f+2+v_2(X_s),\\
 c_E(\ell)&\le {\bf1}_{\ell-1\mid X_s}(1+v_\ell(X_s))
                   \qquad(\ell>2).
 \end{split}
\]
In particular the odd-prime contribution vanishes for
\(\ell>X_s+1\). Summing yields
\[
 \sum_\ell c_E(\ell)\log\ell
 \le (f+2)\log2 + \log X_s+\sum_{\ell\le X_s+1}\log\ell
 \ll b(1+L_b).
\]
Here \(\sum_{\ell\mid X_s}v_\ell(X_s)\log\ell=\log X_s\), and
even the bound \(\sum_{\ell\le X_s+1}\log\ell\le(X_s+1)\log(X_s+1)\)
suffices.

For the remaining part of \(n_s(\ell)\), put
\(\delta_\ell=v_\ell(D_E)/(2f)\) and let \(a_0(\ell)\) be the
maximum local conductor exponent at \(\ell\).
Proposition~\ref{prop:toric-class-denominator}, including the term
\(Wa_0\) and the ceiling to an integer, gives
\[
 n_s(\ell)-c_E(\ell)\ll
 fb\bigl(a_0(\ell)+\delta_\ell+
 {\bf1}_{a_0(\ell)>0}+{\bf1}_{\delta_\ell>0}\bigr).
\]

Indeed \(W=X_s+2K_\kappa\ll fb\), both indicators in the local
ramified allowance are at most one, and the ceiling costs at most one
when the ramified allowance is positive. The conductor terms are
supported on \(S_*\). For the discriminant terms use that \(E/\Q\)
is Galois. At a ramified prime of common ramification index
\(e\ge2\), the local different exponent is at least \(e-1\), hence
\(\delta_\ell\ge(e-1)/e\ge1/2\). Consequently
\[
 \sum_\ell \delta_\ell\log\ell=\log\operatorname{rd}_E,
 \qquad
 \sum_{\ell:\,\delta_\ell>0}\log\ell\le2\log\operatorname{rd}_E .
\]

This proves \(D_T=\sum_\ell n_s(\ell)\log\ell\ll fb(1+L_b)\).
Choose the common comparison constant so that
\(A_T,D_T\le(C/6)fb(\log b+\log x)\).
Writing \(m_{c,w}=m_w+v_w(d_c)\), truncation toward zero gives
\[
 |g_w|\le|m_w|,\qquad
 |v_w(d_c)+g_w|\le |m_{c,w}|+|m_w|.
\]

The two absolute-value estimates proved above now yield
\begin{equation}\label{eq:assigned-weighted-deletions}
 \begin{gathered}
  \sum_w\omega_w |g_w|\le B_T,\qquad
  \sum_w\omega_w|v_w(d_c)+g_w|\le B_T\quad(c=0,1),\\
  B_T=Cfb(\log b+\log x).
 \end{gathered}
\end{equation}

The proof pays for the two branches by summing over their
respective subsets of places. Both use the shift from the
base minimal polynomial. \label{cor:fitted-square-ratios-6}

\end{proof}

\begin{proposition}[Uniform valuation bound for the fitted square]\label{prop:fitted-square-valuation}
For the assigned tuples of Proposition~\ref{prop:fitted-square-ratio-identity},
with the common index cutoff, use the bottom and toric bounds of
Propositions~\ref{prop:fitted-square-bottom-index}
and~\ref{prop:fitted-toric-assigned-bounds}. For the fitted entry, the estimate \eqref{eq:fitted-square-valuation} is
\[
 \begin{aligned}v_w(r_{t,ij})\le{}&|v_w(A)|+|v_w(A')|+4\sum_{s=1}^{H}a_s(w)+5e_w\sum_{s=1}^{H}n_s(\ell)\\&+2He_w/\ell^2+6\delta_w+(p+q)B_w^{\rm step}.\end{aligned}
\]
 The same bound with zero
on the left holds at the base tuple. One threshold works over
the square indices, coefficient places and their assigned tuples.
\end{proposition}

\begin{proof}
\proofstep{Count the positive and negative multiplicities of the index intervals.}

Use the bottom comparison of Proposition~\ref{prop:fitted-square-bottom-index}
and the two-sided bound of Proposition~\ref{prop:fitted-toric-assigned-bounds}.

Fix \(t=t(w)\) with extra level prime to \(\ell\). Put
\(V_s=v_w(\rho_t(s))\) for this same tuple at every index. Before
estimating any of the sums, \eqref{eq:fitted-square-tower} gives
\[
 v_w(r_{t,ij})=v_w(\gamma_t)+\sum_{k=1}^{p+q}v_w(u_k)
       +\sum_{s=e_0+1}^{e+q}V_s-\sum_{s=e+1}^{e+q}V_s .
\]

The unit \(\zeta\) has zero valuation here. In
\eqref{eq:fitted-square-bottom}, reflection sends the indices
\(1-e_0,\ldots,0\) bijectively to \(1,\ldots,e_0\) without
changing \(V_s\). Thus the total signed contribution at the positive
indices is exactly
\[
\begin{aligned}
 \sum_{s=1}^{h_0}V_s+\sum_{s=1}^{e_0}V_s
 +\sum_{s=e_0+1}^{e+q}V_s-\sum_{s=e+1}^{e+q}V_s
 &=\sum_{s=1}^{h_0}V_s+\sum_{s=1}^{e}V_s .
\end{aligned}
\]

Indeed \(e_0+p=e+q\). The upper endpoints lie in
\([1,H]\) since \(e_0\le h_0\) and \(H=h_0+p+q\le N\).
The net multiplicity is at most two and is nonnegative,
including when either set of removals is empty. Hence the upper
bound in \eqref{eq:fitted-toric-two-sides} alone bounds the last
display by
\[
 2\sum_{s=1}^H a_s(w)+2H e_w/\ell^2 .
\]

Each step costs at most \(B_w^{\rm step}\) at this place. Adding
the two base norm costs and \(6\delta_w\) from
\eqref{eq:fitted-square-bottom}, and using \(a_s(w),n_s(\ell)\ge0\),
gives the estimate used for the full grid.
\begin{align}\label{eq:fitted-square-valuation}
 v_w(r_{t,ij})\le{}&
   |v_w(A)|+|v_w(A')|+4\sum_{s=1}^{H}a_s(w)
       +5e_w\sum_{s=1}^{H}n_s(\ell)\\
 &\quad+2H e_w/\ell^2+6\delta_w+(p+q)B_w^{\rm step}.\notag
\end{align}

The bound is also valid in the base case, with zero left side.
The defect, denominators and toric shifts are those of the
same fitted square. All the statements here use the one
grid cutoff and its eventual size bounds, independent of the
index \(i,j\), the place \(w\) and the chosen tuple at that place.
In particular the hypotheses in the upper toric test are
required at the places with their assigned twists.
\label{cor:fitted-square-ratios-7}
\end{proof}
